\section{Versuchsaufbau und Evaluation}
\label{sec:experimental-setup}

\subsection{Evaluationsprotokoll}
\label{sec:cv}

Alle Kennzahlen stammen aus einer gruppierten Five-Fold Cross-Validation über
die Personen, abgestimmt mit der Partnergruppe an der ZHAW. Sämtliche Fenster
einer Person liegen entweder im Training oder im Test, jede der 32 Personen
ist genau einmal Testperson. Random Forest und Netze wurden mit
unterschiedlichen Fold-Partitionen gerechnet; gepaarte Tests zwischen den
Familien sind auf dieser Kohorte daher nicht möglich. Ein früher verwendetes
Leave-One-Subject-Out lieferte auf denselben Daten praktisch identische
Werte.

Glättung, Kalibrierung und Fehleranalyse rechnen auf der Out-of-Fold-Tabelle
der Fensterwahrscheinlichkeiten und erben dieselbe Leakage-Garantie; die
Übergangsmatrix des HMM wird je Runde nur aus den Trainingspersonen
geschätzt.

\subsection{Metriken und statistische Analyse}
\label{sec:metrics}

Primäre Kennzahl ist die Accuracy je Fenster bei fester Schwelle 0,5, weil
die Klassen nahezu ausgeglichen sind und ein Tracker eine Entscheidung
treffen muss. Ergänzend berichten wir ROC-AUC und den F1-Wert der Klasse
Schreiben, jeweils mit ihrer Entscheidungsskala.

Die Streuung zwischen Personen liegt bei drei bis vier Prozentpunkten.
Jeden Vergleich auf denselben Personen prüfen wir deshalb mit dem
Wilcoxon-Vorzeichen-Rang-Test \citep{wilcoxon1945} auf den personenweisen
Differenzen und nennen, bei wie vielen Personen eine Konfiguration vorn
liegt. Ein Gewinn unter einem Prozentpunkt ohne $p < 0{,}05$ gilt als
Rauschen. Für die Netze mitteln wir drei Seeds, weil zwei GPU-Läufe mit
identischem Seed um bis zu 1,6 Prozentpunkte auseinanderliegen können; bei
32 Personen liegen für die Netze nur Seed-Aggregate vor, ihr Vergleich
untereinander ist dort nominal.

\paragraph{Modellauswahl auf derselben Kohorte.} Architekturen und
Hyperparameter wurden anhand derselben Cross-Validation ausgewählt, auf der
wir berichten. Personen kreuzen nie zwischen Training und Test, aber die Werte der Netze
sind Selektionssieger und damit nach oben verzerrt; eine Nested
Cross-Validation war mit dem Rechenbudget nicht darstellbar. Die Sieger der
Hyperparametersuche verloren beim Nachrechnen mit drei Seeds ein bis zwei
Prozentpunkte gegenüber ihrem besten Einzellauf. Die Hyperparameter des
Random Forest blieben über das Projekt unverändert.

\subsection{Baselines und Ablationsstudien}
\label{sec:ablations}

Der Random Forest auf 88 Merkmalen ist die Referenz. Die Netze auf nativen
5-s-Fenstern vergleichen wir mit ihm bei gleicher Entscheidungslatenz: gegen
seine kausal über 5~s geglätteten 1-s-Vorhersagen und gegen einen Random
Forest mit direkt über 5~s berechneten Merkmalen. Die Ablationen ändern je
eine Variable auf identischen Folds und Seeds: Schwerkraftvektor,
Gyroskop, Dekodiermodus und Fusion. Zwei Vergleiche sind nicht auf eine
Variable beschränkt (50 gegen 100~Hz, Drei-Kanal-Lauf der Netze) und werden
entsprechend eingeordnet. Wirkungslose Eingriffe fasst Abschnitt~\ref{sec:ablation-results} zusammen,
weil sie die Interpretation der Leistungsgrenze tragen. Random Forest und HMM
laufen lokal in unter einer Stunde, die Netze brauchten auf einer
Cloud-GPU rund zehn Stunden.
